
\documentclass[
  reprint,
  amsmath,amssymb,
  aps,
  prd,
  superscriptaddress,
  nofootinbib,
  longbibliography
]{revtex4-2}

% -------------------- PACKAGES --------------------
\usepackage[T1]{fontenc}
\usepackage[utf8]{inputenc}
\usepackage{lmodern}
\usepackage{microtype}
\usepackage{bm}
\usepackage{mathtools} % loads amsmath
\usepackage{amsthm}
\usepackage{siunitx}
\sisetup{uncertainty-mode = separate, per-mode = symbol}
\usepackage{graphicx}
\usepackage[dvipsnames]{xcolor}
\usepackage[hidelinks]{hyperref}
\usepackage{physics}
\usepackage{booktabs}
\usepackage{orcidlink}

% -------------------- MACROS --------------------
\newcommand{\phig}{\varphi}
\newcommand{\Sthree}{S^{3}}
\newcommand{\ellP}{\ell_{\rm P}}
\newcommand{\tP}{t_{\rm P}}
\newcommand{\Sig}{\Sigma}
\newcommand{\Diff}{\mathrm{Diff}}
\newcommand{\Lie}{\mathcal L}
% \dd provided by physics package
\newcommand{\RR}{\mathbb R}
\newcommand{\ZZ}{\mathbb Z}

% -------------------- THEOREM STYLES --------------------
\newtheorem{theorem}{Theorem}
\newtheorem{lemma}{Lemma}
\newtheorem{corollary}{Corollary}
\newtheorem{proposition}{Proposition}
\theoremstyle{definition}
\newtheorem{definition}{Definition}
\theoremstyle{remark}
\newtheorem{remark}{Remark}

\begin{document}

% Manuscript title (classic PRD style; also make PDF bookmarks math-safe)
\title{Noether Identities and Discrete \texorpdfstring{$\varphi$}{phi}-Embedding in Compact \texorpdfstring{$S^{3}$}{S3} Quantum Cosmology}

\author{Andrei Preece\orcidlink{0009-0004-5978-795X}}
\affiliation{Independent Researcher, Berlin, Germany}

\author{Boris Batenin\orcidlink{0009-0000-7027-8583}}
\affiliation{Independent Researcher, Berlin, Germany}

\date{\today}
\begin{abstract}
We present a topology-explicit and Noether-consistent variational framework for quantum cosmology on
a spacetime $M=\RR\times \Sig$ with compact spatial hypersurfaces $\Sig$, and we make explicit the
topological sector in which the universal cover is spherical, $\widetilde{\Sig}\simeq \Sthree$.
General covariance implies a Noether-II identity structure (Bianchi-type
identities) rather than genuine global conserved charges. We then introduce a discrete $\phig$-embedding level
$n\in\ZZ$ (or a time-dependent embedding $N(t)$) within the action and state precise conditions under which an
on-shell conserved current $\nabla_\mu J^\mu_{(n)}=0$ exists. The discrete choice
$N(t)=\lfloor \log_{\phig}(ct/\ell_P)\rfloor$ is shown to require explicit junction conditions at transition times
to prevent distributional sources in the Euler--Lagrange equations. Alternatively, a smooth embedding variable
$\nu(t)=\log_{\phig}(ct/\ell_P)$ yields a fully differentiable action, with discretization emerging at the level
of spectra. We derive the boundary/topology corollary: for compact $\Sig$ with $\partial\Sig=\varnothing$,
diffeomorphism-Noether charges reduce to boundary integrals and vanish globally (in particular for
$\Sig\cong \Sthree$). Finally, we provide minimal,
falsifiable observational discriminants in CMB anisotropies and stochastic gravitational-wave backgrounds
sensitive to compact topology and discrete mode structure, and we outline a likelihood-based verification pipeline.
\end{abstract}

\maketitle

\section{Motivation and scope}
In generally covariant theories such as General Relativity (GR), the relevant symmetry is local:
invariance under diffeomorphisms $\Diff(M)$. Noether's second theorem~\cite{Noether1918} therefore applies: the symmetry implies
differential identities among the Euler--Lagrange (EL) expressions rather than independent global conservation
laws of Noether-I type. This distinction becomes operationally critical in cosmology, where global notions such as
``total energy of the Universe'' depend sensitively on the existence of boundaries and the global topology of
spatial hypersurfaces $\Sig$.

The aim of this manuscript is twofold:
(i) to state the Noether-II content of GR in a topology-explicit way (boundary charges and their vanishing on
compact boundaryless $\Sig$ with $\partial\Sig=\varnothing$, including the spherical sector $\widetilde{\Sig}\simeq\Sthree$), and
(ii) to formulate a mathematically consistent implementation of a discrete $\phig$-embedding label $n$ (or $N(t)$)
inside a variational principle, including the required junction conditions for piecewise-smooth embeddings.
We then connect the framework to observational discriminants of compact topology: discrete Laplacian spectra,
low-$\ell$ CMB structure, and infrared properties of the stochastic gravitational-wave background.

\subsection{Invariant fixing of the golden ratio scale factor}
\label{sec:phi_invariant}
The discrete embedding label $n$ is not introduced as a tunable convention, but as an invariant bookkeeping of
\emph{scale composition}: coarse-graining steps must add without ambiguity, and the resulting representation of
integers must be unique under refinement. The mathematically sharp form of this requirement is the
\emph{Zeckendorf property}: every positive integer admits a unique decomposition into a sum of \emph{non-adjacent}
elements of a strictly increasing ``scale sequence.'' In the classical case this sequence is the Fibonacci ladder,
and uniqueness is guaranteed by the Zeckendorf--Lekkerkerker theorem~\cite{Lekkerkerker1951, Keller1972}.

Physically, this ``unique decomposition'' is the minimal invariance needed for a well-defined refinement
semigroup: successive coarse-grainings must compose associatively and yield a path-independent index of scale.
In quantum cosmology this provides an invariant coarse-graining coordinate (a discrete RG-like clock) rather than
an arbitrary choice of base.

To avoid confusion with the large ``golden-ratio cosmology literature (often based on post-hoc pattern matching),
we stress that $\phig$ does not enter here as a fit parameter or an empirical coincidence. Instead, $\phig$ is a
\emph{rigidity consequence} of imposing a minimal invariance requirement on scale composition (unique, path-independent
refinement in a discrete semigroup) \emph{inside} the variational principle. The empirical content of the framework is
therefore not ``$\phig$ appears in data, but rather the falsifiable prediction that compact-topology observables must
be computed with discrete mode sums and (if present) $\phig$-indexed features, and then tested against likelihoods
(Secs.~\ref{sec:likelihood-protocol} and~\ref{sec:sgwb}).

A key rigidity statement (a weak converse) shows that this uniqueness requirement is not merely consistent with the
Fibonacci ladder but essentially characterizes it among monotone sequences~\cite{Chang2020WeakConverse}. Hence the
discrete scale ladder is fixed (up to initial normalization) by invariance/uniqueness alone. The associated
\emph{scale multiplier} is then forced as the unique positive fixed point of the induced recursion:
\begin{lemma}[Golden ratio as the invariant multiplier of the Zeckendorf ladder]
\label{lem:phi_multiplier}
Let $\{F_n\}$ be the Fibonacci sequence defined by $F_{n+2}=F_{n+1}+F_n$ with $F_0=F_1=1$. Then the ratio
$r_n=F_{n+1}/F_n$ converges to the unique positive root of $r^2=r+1$, namely
\begin{equation}
\lim_{n\to\infty}\frac{F_{n+1}}{F_n}=\phig\equiv \frac{1+\sqrt{5}}{2}.
\end{equation}
\end{lemma}

In the present framework, $n$ labels this invariant discrete ladder of scale composition; therefore $\phig$ enters
as the uniquely selected multiplier of the ladder, and the smooth embedding $\nu(t)=\log_{\phig}(ct/\ell_P)$ is the
coordinate-free continuum interpolation of that invariant index.

\section{Geometric setup and variational calculus}
\subsection{Spacetime and topology}
\begin{definition}[Spacetime decomposition]
Let $M=\RR\times \Sig$, where $\Sig$ is a compact orientable 3-manifold \emph{without boundary}
($\partial\Sig=\varnothing$). The Noether-II/topology statements below depend only on this boundarylessness.
\end{definition}

\begin{proposition}[Spherical cover as an invariantly characterized sector]
\label{prop:spherical_space_form}
If $\Sig$ admits a (complete) Riemannian metric of constant \emph{positive} sectional curvature, then
$\Sig$ is a spherical space form: there exists a finite group $\Gamma\subset \mathrm{Isom}(\Sthree)$ acting freely
such that
\begin{equation}
\Sig\simeq \Sthree/\Gamma,
\end{equation}
and hence the universal cover is $\widetilde{\Sig}\simeq \Sthree$.\cite{WolfConstantCurvature,Scott3Manifolds}
\end{proposition}

In the observational part of the paper, this spherical-cover sector is singled out by its invariant spectral
fingerprints (discrete Laplacian eigenvalues and characteristic degeneracies), which are directly testable in
low-$\ell$ CMB statistics and in infrared features of $\Omega_{\rm GW}(f)$.

\subsection{Action and EL expressions}
Consider a generally covariant action
\begin{equation}
S[g,\psi]=\int_{M}\dd^4x\,\sqrt{-g}\,\mathcal L(g,\psi,\nabla\psi,\nabla\nabla\psi,\dots),
\label{eq:action_general}
\end{equation}
where $\psi$ denotes matter or effective cosmological degrees of freedom. We define EL expressions as
\begin{equation}
E^{\mu\nu}\equiv \frac{1}{\sqrt{-g}}\frac{\delta S}{\delta g_{\mu\nu}},
\qquad
E_{\psi}\equiv \frac{1}{\sqrt{-g}}\frac{\delta S}{\delta \psi}.
\label{eq:EL_defs}
\end{equation}

\section{Noether II for diffeomorphism invariance}
\subsection{Local symmetry and identities}
\begin{theorem}[Noether II for diffeomorphism invariance]
\label{thm:noetherII_diff}
Let $S[g,\psi]$ be invariant under infinitesimal diffeomorphisms generated by an arbitrary smooth vector field
$\xi^\mu(x)$:
\begin{equation}
\delta_\xi g_{\mu\nu}=\nabla_\mu\xi_\nu+\nabla_\nu\xi_\mu,
\qquad
\delta_\xi \psi=\Lie_\xi \psi.
\label{eq:diff_variations}
\end{equation}
Then the EL expressions satisfy differential identities (Noether-II identities) of the schematic form
\begin{equation}
\nabla_\mu E^{\mu\nu} \;+\; \mathcal K^\nu(E_\psi,\nabla E_\psi,\dots)\;\equiv\;0,
\label{eq:noetherII_identity}
\end{equation}
where $\mathcal K^\nu$ is a covariant differential operator determined by the tensor character of $\psi$.
\end{theorem}

\begin{proof}
Diffeomorphism invariance depends on arbitrary functions $\xi^\mu(x)$, hence it falls under Noether's second theorem~\cite{Noether1918}.
Varying $S$ under \eqref{eq:diff_variations} yields the schematic identity
\begin{align}
0=\delta_\xi S
&=\int_M \dd^4x\,\sqrt{-g}\Big(E^{\mu\nu}\delta_\xi g_{\mu\nu}+E_\psi\,\delta_\xi\psi\Big)
\nonumber\\
&\quad+\int_M \dd^4x\,\sqrt{-g}\,\nabla_\mu(\cdots).
\end{align}
Integrating by parts transfers derivatives from $\xi^\mu$ to $E^{\mu\nu}$ and $E_\psi$. Since $\xi^\mu(x)$ is
arbitrary, the coefficient of $\xi^\mu$ must vanish identically, giving \eqref{eq:noetherII_identity}.
\end{proof}

\subsection{Boundary charges, compact \texorpdfstring{$\Sig$}{Sigma}, and the meaning of ``conservation''}
In GR, diffeomorphism invariance is a local gauge symmetry. Consequently, the standard Noether-I expectation
(``rigid symmetry $\Rightarrow$ a global conserved charge'') does not apply in a straightforward way.
Noether's second theorem instead yields identity structure among the Euler--Lagrange expressions, and the
corresponding ``conservation laws'' are \,improper\, in the classic sense: on-shell they reduce to total divergences
and/or boundary terms rather than providing a unique local, gauge-invariant gravitational energy density
\cite{Noether1918,BradingBrown2003}.

A precise formulation is given by the covariant phase space approach \cite{LeeWald1990,IyerWald1994,WaldBook}.
Treat the Lagrangian as an $n$-form $L$ and write
\begin{equation}
\delta L \;=\; E(\phi)\,\delta\phi \; +\; d\Theta(\phi,\delta\phi),
\end{equation}
where $E(\phi)=0$ are the field equations and $\Theta$ is the symplectic potential. For a vector field $\xi$
generating an infinitesimal diffeomorphism, define the Noether current $(n\! -\! 1)$-form
\begin{equation}
J_{\xi} \;\equiv\; \Theta(\phi,\Lie_{\xi}\phi)\; -\; \xi\cdot L.
\end{equation}
Diffeomorphism invariance implies that $J_{\xi}$ is closed on-shell, $dJ_{\xi}\approx 0$, and, crucially, on-shell
\emph{exact}:
\begin{equation}
J_{\xi} \;\approx\; dQ_{\xi},
\label{eq:noether_current_exact}
\end{equation}
with $Q_{\xi}$ an $(n\! -\! 2)$-form (Noether charge form). This is the clean statement that covariant charges in GR
are generically boundary charges.

At the level of canonical generators, one may write schematically
\begin{equation}
H_{\xi} \;=\; \int_{\partial\Sig} \bigl(Q_{\xi} - \xi\cdot B\bigr),
\label{eq:hamiltonian_boundary}
\end{equation}
where $B$ is a boundary counterterm fixed by the admissible boundary conditions \cite{IyerWald1994,WaldBook}.
Equation \eqref{eq:hamiltonian_boundary} makes the compact-sector corollary immediate.

\begin{corollary}[Compact boundaryless $\Sig$]
\label{cor:S3_charge_zero}
If $\Sig$ is compact and boundaryless, $\partial\Sig=\varnothing$, then diffeomorphism-Noether charges reduce to
boundary integrals and vanish globally:
\begin{equation}
\partial\Sig=\varnothing \qquad\Longrightarrow\qquad H_{\xi}=0.
\end{equation}
\noindent(For the corresponding $\xi$ and standard boundary conditions.)
In particular, for $\Sig\cong\Sthree$ (or more generally $\Sig\simeq\Sthree/\Gamma$), there is no global Noether-I
``energy of the Universe''; the invariant content is the Noether-II identity structure in
Theorem~\ref{thm:noetherII_diff}.
\end{corollary}

\begin{remark}[Junction consistency under discrete refinement]
When refinement is implemented piecewise (e.g.\ via the floor embedding $N(t)$ below), Noether-consistency requires
that matching data do not introduce spurious distributional sources. In the covariant language this means that the
exactness property \eqref{eq:noether_current_exact} is not spoiled by junction contributions beyond admissible
boundary counterterms. This is the conceptual origin of the explicit junction conditions derived in
Section~IV.
\end{remark}

\section{\texorpdfstring{Discrete $\phig$-embedding: smooth vs floor and junction conditions}{Discrete phi-embedding: smooth vs floor and junction conditions}}
\subsection{Embedding variables}
Define the smooth embedding variable
\begin{equation}
\nu(t)=\log_{\phig}\!\left(\frac{c\,t}{\ell_P}\right),
\label{eq:nu_def}
\end{equation}
and the discrete embedding level
\begin{equation}
N(t)=\lfloor \nu(t)\rfloor,\qquad N\in\ZZ.
\label{eq:N_floor}
\end{equation}

\subsection{\texorpdfstring{Prototype $\phig$-dependent matter sector}{Prototype phi-dependent matter sector}}
As a minimal prototype, consider a complex scalar $\psi$ with a $\phig$-dependent potential:
\begin{align}
\mathcal L
&=\frac12 g^{\mu\nu}\partial_\mu\psi^*\partial_\nu\psi - V(|\psi|;N),
\nonumber\\
V(|\psi|;N)
&=\frac12 m^2|\psi|^2\,\phig^{-2N}.
\label{eq:prototype_L}
\end{align}
More generally, one may encode the embedding level as a conformal scaling
\begin{equation}
g_{\mu\nu}(x,N)=\phig^{-2N}\,\tilde g_{\mu\nu}(x),
\label{eq:metric_scaling}
\end{equation}
provided the physical (nondegenerate) metric entering observables is defined so that degeneracy is avoided.
Concretely, one may treat \(\tilde g_{\mu\nu}\) as the physical metric and interpret \eqref{eq:metric_scaling}
as an \emph{auxiliary} scale bookkeeping that is absorbed into the definition of dimensionless couplings and
potentials; then all observables are built from dimensionless ratios and remain well-defined in the smooth limit.

\subsection{Two consistent formulations}
\paragraph{(I) Smooth embedding.}
\begin{equation}
S=\int \dd^4x\,\sqrt{-g}\,\mathcal L(\phi,\nabla\phi;\nu(t)).
\label{eq:action_smooth_nu}
\end{equation}

\paragraph{(II) Discrete floor embedding.}
\begin{equation}
S=\int \dd^4x\,\sqrt{-g}\,\mathcal L(\phi,\nabla\phi;N(t)).
\label{eq:action_floorN}
\end{equation}

\subsection{Junction term from the variational principle}
Splitting into intervals $I_k=(t_k,t_{k+1})$ and integrating by parts yields
\begin{align}
\delta S
&=\sum_k \int_{I_k}\dd^4x\,\sqrt{-g}\,\Big(E_\phi\,\delta\phi\Big)
\nonumber\\
&\quad+\sum_k \int_{\Sig_{t_k}}\dd^3x\,\left[\Pi_\phi\,\delta\phi\right]_{t_k^-}^{t_k^+},
\label{eq:deltaS_junction}
\end{align}
with
\begin{equation}
\Pi_\phi\equiv \frac{\partial(\sqrt{-g}\mathcal L)}{\partial(\partial_t\phi)}.
\label{eq:Pi_def}
\end{equation}

\begin{theorem}[Junction consistency for distribution-free EL equations]
\label{thm:junction_condition}
In the discrete formulation \eqref{eq:action_floorN}, the EL equations are distribution-free iff
\begin{equation}
\int_{\Sig_{t_k}}\dd^3x\, \Big[\Pi_\phi\,\delta\phi\Big]^{t_k^+}_{t_k^-}=0
\qquad \text{for all }k.
\label{eq:junction_condition}
\end{equation}
A sufficient standard choice is continuity of canonical data:
\begin{equation}
\phi(t_k^-)=\phi(t_k^+),\qquad \Pi_\phi(t_k^-)=\Pi_\phi(t_k^+).
\label{eq:canonical_continuity}
\end{equation}
\end{theorem}

\begin{proof}
\eqref{eq:deltaS_junction} is exact. Stationarity requires both the bulk and junction contributions to vanish for all
admissible variations, giving \eqref{eq:junction_condition}.
\end{proof}

\subsection{Implication for conserved currents}
If on each $I_k$ there is a current $J^\mu_{(N)}$ with $\nabla_\mu J^\mu_{(N)}=0$ on-shell, then global conservation
holds iff
\begin{equation}
\sum_k \int_{\Sig_{t_k}}\dd^3x\,\sqrt{h}\, n_\mu J^\mu_{(N)}\Big|^{t_k^+}_{t_k^-}=0.
\label{eq:current_junction_flux}
\end{equation}


\subsection{Explicit generalized action and scaling symmetry (comparison to the two preprints)}
\label{sec:explicit-generalized-noether}

For completeness, and to make the present manuscript self-contained, we restate the
core structures of the two companion preprints on a ``generalized Noether theorem''~\cite{PreeceBatenin2025QCNoether,PreeceBatenin2025MFNoether}
in the $\mathrm{FBS^3R}$ setting (quantum-cosmology and multifractal versions) in a unified notation.

\paragraph{(A) Fixed embedding level $n$ (quantum-cosmology sketch).}
The preprint model starts from a $\phig$-indexed family of length scales
\begin{equation}
R_n=\ellP\,\phig^{\,n},
\end{equation}
and an action for a complex scalar ``wave function of the Universe'' $\psi$ on a background metric $g_{\mu\nu}$,
\begin{equation}
S_n[\psi,g]=\int d^4x\,\sqrt{-g}\;\mathcal{L}(\psi,\partial_\mu\psi;\phig^{\,n}),
\label{eq:Sn-fixedn}
\end{equation}
with Lagrangian density of the form
\begin{align}
\mathcal{L}
&=\frac12\, g^{\mu\nu}\,\partial_\mu\psi^\ast\,\partial_\nu\psi - V(|\psi|,\phig^{\,n}),
\nonumber\\
V(|\psi|,\phig^{\,n})
&=\frac12\,m^2|\psi|^2\,\phig^{-2n}.
\label{eq:L-fixedn}
\end{align}
The generalized scaling transformation is taken as
\begin{equation}
x^\mu \mapsto \phig\,x^\mu,
\label{eq:scale-fixedn}
\end{equation}
and the associated Noether current (in the standard variational form) is
\begin{equation}
J^\mu_{(n)}=
\frac{\partial\mathcal{L}}{\partial(\partial_\mu\psi)}\,\delta\psi
-
\Big(
\frac{\partial\mathcal{L}}{\partial(\partial_\mu\psi)}\,\partial_\nu\psi
-
\mathcal{L}\,g^\mu{}_\nu
\Big)\delta x^\nu .
\label{eq:noether-current-fixedn}
\end{equation}
On shell this satisfies the covariant conservation law
\begin{equation}
\nabla_\mu J^\mu_{(n)}\approx 0,
\label{eq:current-conservation-fixedn}
\end{equation}
where $\approx$ denotes equality modulo the Euler--Lagrange equations.

\paragraph{(B) Multifractal action with a time-dependent embedding level $N(t)$.}
The multifractal preprint specifies a metric scaling
\begin{equation}
ds^2=g_{\mu\nu}(x,n)\,dx^\mu dx^\nu,
\qquad
g_{\mu\nu}(x,n)=\phig^{-2n}\,\tilde g_{\mu\nu}(x),
\label{eq:metric-multifractal}
\end{equation}
and a variational principle
\begin{equation}
S=\int d^4x\,\sqrt{-g}\;\mathcal{L}(\psi,\partial_\mu\psi,n).
\label{eq:action-multifractal}
\end{equation}
If $n=n(\psi)$ is treated as an additional dependent variable, the Euler--Lagrange equation takes the modified form
\begin{equation}
\frac{\partial\mathcal{L}}{\partial\psi}-\nabla_\mu\!\Big(\frac{\partial\mathcal{L}}{\partial(\partial_\mu\psi)}\Big)
+\frac{\partial\mathcal{L}}{\partial n}\,\frac{\partial n}{\partial\psi}=0.
\label{eq:EL-multifractal}
\end{equation}
The embedding level is then promoted to a time-dependent quantity
\begin{equation}
\begin{split}
N(t)
&=\Big\lfloor\log_{\frac{1+\sqrt5}{2}}\!\Big(\frac{ct}{\ellP}\Big)\Big\rfloor
 =\Big\lfloor\log_{\phig}\!\Big(\frac{ct}{\ellP}\Big)\Big\rfloor,
\end{split}
\label{eq:Nt-floor}
\end{equation}
and the action is written as
\begin{equation}
\begin{split}
S_{\mathrm{FBS^3R}}
&=\int d^4x\,\sqrt{-g_{N(t)}}\;
\mathcal{L}_{N(t)}\!\Big(\psi,\partial_\mu\psi,\phig^{N(t)}\Big),
\end{split}
\label{eq:action-Nt}
\end{equation}
with
\begin{align}
\mathcal{L}_{n}
&=\frac12\,g^{\mu\nu}\partial_\mu\psi^\ast\partial_\nu\psi
-V\!\big(|\psi|,\phig^{n}\big),
\nonumber\\
V\!\big(|\psi|,\phig^{n}\big)
&=\frac12\,m^2|\psi|^2\,\phig^{-4n}.
\label{eq:L-Nt}
\end{align}
The preprint further proposes an illustrative two-regime ansatz for the wave function,
\begin{equation}
\psi(t)=
\begin{cases}
A\left(\dfrac{ct}{\ellP}\right)^{-3/2}\exp\!\Big(i\,\phig^{N(t)}\,\dfrac{t}{\tP}\Big),
& t<t_{\rm rec},
\\[0.6em]
B\,\exp\!\big(-\tfrac32 H(t-t_{\rm rec})\big)\,\exp(iHt), & t\ge t_{\rm rec},
\end{cases}
\label{eq:psi-ansatz}
\end{equation}

In the ``large-$N$'' continuum approximation used for intuition in the multifractal sketch, one writes
\begin{align}
\exp\!\Big(i\,\phig^{N(t)}\frac{t}{\tP}\Big)&\;\sim\;\exp\!\Big(i\,\frac{ct}{\ellP}\,\ln t\Big),
\label{eq:phase-largeN-approx}
\\
g_{\mu\nu}(x,N(t))=\phig^{-2N(t)}\tilde g_{\mu\nu}(x)
&\;\sim\;\left(\frac{ct}{\ellP}\right)^{-2}\tilde g_{\mu\nu}(x),
\label{eq:metric-largeN-approx}
\\
V\big(|\psi|,\phig^{N(t)}\big)=\frac12 m^2|\psi|^2\phig^{-4N(t)}
&\;\sim\;\frac12 m^2|\psi|^2\left(\frac{ct}{\ellP}\right)^{-4},
\label{eq:V-largeN-approx}
\end{align}
which should be treated as a heuristic scaling limit, while the discrete $N(t)$ formulation requires the
junction consistency already stated in Sec.~IV.
where $t_{\rm rec}$ is a reference transition time and $H$ is the Hubble scale.

The multifractal scaling symmetry is specified as
\begin{equation}
\begin{split}
x^\mu \mapsto \phig\,x^\mu,\qquad
N(t)\mapsto N(t)+1,\qquad
\psi\mapsto \psi/\phig^{3/2},
\end{split}
\label{eq:scale-multifractal}
\end{equation}
leading to the current
\begin{align}
\mathcal{J}^\mu_{n}
&=\frac{\partial \mathcal{L}_{n}}{\partial(\partial_\mu\psi)}\,\delta\psi
-
\Big(
\frac{\partial \mathcal{L}_{n}}{\partial(\partial_\mu\psi)}\,\partial_\nu\psi
-
\mathcal{L}_{n}\,g^\mu{}_\nu
\Big)\delta x^\nu,
\nonumber\\
\nabla_\mu \mathcal{J}^\mu_{n}&=0.
\label{eq:noether-current-multifractal}
\end{align}
and a corresponding conserved quantity (for a chosen time foliation)
\begin{equation}
E_n=\int d^3x\,\sqrt{-g_n}\;\mathcal{J}^0_{n}.
\label{eq:En}
\end{equation}

\paragraph{What is retained, and what is strengthened here.}
Equations \eqref{eq:Sn-fixedn}--\eqref{eq:En} are the precise elements of the two preprints that are
\emph{needed} for the unified PRD manuscript: (i) the explicit symmetry map
\eqref{eq:scale-fixedn} and/or \eqref{eq:scale-multifractal}, (ii) the explicit current formulas
\eqref{eq:noether-current-fixedn} and \eqref{eq:noether-current-multifractal}, and (iii) the exact
definition of the time-dependent embedding \eqref{eq:Nt-floor}.
The present manuscript goes beyond the sketches by (a) clarifying the on-shell nature of the current
identities, (b) making junction conditions for floor-discretized $N(t)$ explicit (Theorem~\ref{thm:junction_condition}),
and (c) separating \emph{predictions} from \emph{tests} via a likelihood protocol (Sec.~\ref{sec:likelihood-protocol}).

\begin{remark}[Covariance in the presence of a discrete embedding]
The underlying gravitational sector remains diffeomorphism-invariant. In the \emph{smooth} formulation one may treat
$\nu(t)=\log_{\phig}(ct/\ellP)$ as an (effective) scalar embedding variable. In the \emph{floor} formulation, $N(t)$ is
piecewise-constant with a chosen foliation; adding the junction term and imposing the matching conditions ensure that the
Euler--Lagrange system is distribution-free, so the Noether-II identity structure holds on each interval and is globally
consistent across the transition hypersurfaces.
\end{remark}


\section{Topology and CMB observables}

\subsection{Invariant routes to the spherical sector}
\label{sec:invariant_routes_s3}
Diffeomorphism invariance in GR constrains the \emph{local} variational structure via a Noether-II identity
(Bianchi-type relations), but it does not by itself fix the \emph{global} topology of $\Sig$.
To make the spherical sector mathematically unavoidable, one must add genuinely
topological/geometric invariant conditions. We record two standard levels of ``inevitability.''

\paragraph{Hard inevitability (strict): $\Sig\cong \Sthree$.}
\begin{theorem}[Hard spherical inevitability]
\label{thm:hard_s3_main}
Let $\Sig$ be a closed 3-manifold, i.e.\ compact and without boundary, $\partial\Sig=\varnothing$.
If $\Sig$ is simply connected,
\begin{equation}
\pi_1(\Sig)=0,
\end{equation}
then $\Sig$ is homeomorphic to the 3-sphere,
\begin{equation}
\Sig \cong \Sthree.
\end{equation}
\end{theorem}
This is the Poincar\'e/Geometrization criterion in dimension three~\cite{Perelman2002Entropy,Perelman2003Surgery,Perelman2003Extinction,MorganTian2007RicciFlowPoincare,KleinerLott2008NotesPerelman}.

\paragraph{Soft inevitability (minimal spherical): $\widetilde{\Sig}\cong \Sthree$.}
\begin{corollary}[Soft spherical inevitability: spherical space forms]
\label{cor:soft_s3_main}
Let $\Sig$ be closed, $\partial\Sig=\varnothing$. If $\Sig$ admits a Riemannian metric with strictly positive
Ricci curvature,
\begin{equation}
\mathrm{Ric}>0,
\end{equation}
then $\Sig$ is a spherical space form,
\begin{equation}
\Sig\cong \Sthree/\Gamma,\qquad \widetilde{\Sig}\cong \Sthree,
\end{equation}
where $\Gamma$ is a finite group acting freely by isometries~\cite{Hamilton1982PositiveRicci,Scott3Manifolds,WolfConstantCurvature}.
\end{corollary}
In our observational pipeline we therefore treat $\Sig\cong \Sthree/\Gamma$ as the relevant compact sector and compare
candidate $\Gamma$ via likelihood/evidence (Appendix~\ref{app:hard_soft_s3}).

\subsection{Current observational status of compact topology searches}
\label{sec:topology_status}
Direct topology searches in the CMB target the existence of repeated patterns such as ``matched circles''
in the temperature/polarization sky maps~\cite{CornishSpergelStarkman1998}. Current analyses, including the
Planck topology study, find no statistically significant detection and therefore provide lower bounds on
the size (injectivity radius) of possible fundamental domains~\cite{Planck2015Topology}. Recent dedicated
limits on spherical lens spaces $L(p,q)$ from circle searches refine these bounds within the spherical sector
$\Sthree/\Gamma$~\cite{Saha2024LensSpaces}, and forecasts quantify the reach of future searches~\cite{FutureTopologySearches}.

\subsection{Curvature radius and mode discreteness}
In a closed FRW slicing introduce curvature radius $R_K$; at present
\begin{equation}
R_{K,0}=\frac{c}{H_0\sqrt{|\Omega_k|}},\qquad \Omega_k<0,
\label{eq:RK0_Omegak}
\end{equation}
On $\Sthree$ scalar harmonics satisfy
\begin{equation}
\Delta_{\Sthree}Y_{n\ell m}=-\frac{n(n+2)}{R_K^2}\,Y_{n\ell m},
\qquad n=0,1,2,\dots.
\label{eq:S3_scalar_eigs}
\end{equation}
Define $k_n=(n+1)/R_K$.

\subsection{A concrete toy numerical example: discrete \texorpdfstring{$S^{3}$}{S3} mode sum}
\label{sec:toy_example_modes}
To make the discrete mode structure fully explicit (independently of any Boltzmann-code implementation),
consider the scalar Laplacian spectrum \eqref{eq:S3_scalar_eigs}. The associated ``wavenumber'' is
$k_n=(n+1)/R_K$ and the scalar harmonic degeneracy on $\Sthree$ is $(n+1)^2$.
Table~\ref{tab:toy_modes} lists the first modes for an illustrative curvature radius $R_K=\SI{50}{Gpc}$,
and Fig.~\ref{fig:toy_S3_modes} visualizes the discretization (a uniform spacing $\Delta k=1/R_K$).
This is a minimal, checkable numerical instance of the replacement $k\to k_n$ used in
\eqref{eq:Cl_discrete_sum}.

\begin{table}[t]
\centering
\caption{Toy $\Sthree$ scalar modes for $R_K=\SI{50}{Gpc}$. Here $k_n=(n+1)/R_K$ and the scalar degeneracy is $(n+1)^2$.}
\label{tab:toy_modes}
\begin{tabular}{@{}cccc@{}}
\toprule
$n$ & $k_n\;[\mathrm{Gpc}^{-1}]$ & $\Delta k\;[\mathrm{Gpc}^{-1}]$ & degeneracy $(n+1)^2$\\
\midrule
0  & 0.02 & 0.02 & 1  \\
1  & 0.04 & 0.02 & 4  \\
2  & 0.06 & 0.02 & 9  \\
3  & 0.08 & 0.02 & 16 \\
4  & 0.10 & 0.02 & 25 \\
5  & 0.12 & 0.02 & 36 \\
6  & 0.14 & 0.02 & 49 \\
7  & 0.16 & 0.02 & 64 \\
8  & 0.18 & 0.02 & 81 \\
9  & 0.20 & 0.02 & 100\\
10 & 0.22 & 0.02 & 121\\
\bottomrule
\end{tabular}
\end{table}

\begin{figure}[t]
\centering
\includegraphics[width=\columnwidth]{figures/toy_S3_modes.pdf}
\caption{Illustrative discretization of scalar modes on $\Sthree$ for $R_K=\SI{50}{Gpc}$: $k_n=(n+1)/R_K$.}
\label{fig:toy_S3_modes}
\end{figure}


\subsection{\texorpdfstring{Low-$\ell$ scale and $C_\ell$ template}{Low-l scale and Cl template}}
With $\ell\sim k\chi_*$, the fundamental scale is
\begin{equation}
\ell_{\min}\sim \frac{\chi_*}{R_K}.
\label{eq:lmin}
\end{equation}
The angular spectrum becomes
\begin{equation}
C_\ell = 4\pi \sum_{n} \mathcal P_{\mathcal R}(k_n)\,\abs{\Delta_\ell(k_n)}^2.
\label{eq:Cl_discrete_sum}
\end{equation}

\subsection{\texorpdfstring{Matched circles for $\Sig\simeq \Sthree/\Gamma$}{Matched circles for Sigma ≃ S3/Gamma}}
If $\Sig\simeq \Sthree/\Gamma$ with nontrivial $\Gamma$, matched-circle correlations~\cite{CornishSpergelStarkman1998,Planck2015Topology} may appear; non-detection yields
constraints on the fundamental domain size.

\subsection{Minimal Planck likelihood pipeline}
\begin{enumerate}
\item Choose $(R_K,\Omega_k)$ and $\mathcal P_{\mathcal R}(k)$.
\item Compute closed-geometry transfer functions and evaluate \eqref{eq:Cl_discrete_sum}.
\item Fit low-$\ell$ likelihood and constrain $R_K$ (and embedding selection parameters).
\item Apply matched-circle constraints if $\Gamma\neq 1$.
\end{enumerate}


\subsection{Likelihood protocol for observational verification}
\label{sec:likelihood-protocol}

The two companion preprints contained \emph{qualitative} ``alignment'' statements with CMB and GW data.
For PRD-grade verification we adopt an explicit likelihood protocol. The minimal requirement is that
the model be implemented as a prediction map $\theta\mapsto \{C_\ell,\Omega_{\rm GW}(f),\dots\}$ and
tested against public likelihoods.

\paragraph{CMB.}
Let $\mathbf{D}$ denote the binned Planck~\cite{Planck2018Cosmo,Planck2018Lik} data vector (e.g.\ TT,TE,EE) and $\mathbf{M}(\theta)$ the model
prediction obtained from a Boltzmann code (CLASS/CAMB)~\cite{Blas2011CLASS,Lewis2000CAMB} in a closed geometry, with the replacement $k\to k_n=(n+1)/R_K$
in the mode sums \eqref{eq:Cl_discrete_sum} and a parameterization of the primordial spectrum
$\mathcal{P}_{\mathcal{R}}(k)$ that allows for the $\phig$-indexed features.
The (Gaussian-approximate) log-likelihood reads
\begin{equation}
-2\ln\mathcal{L}_{\rm CMB}(\theta)=
\big(\mathbf{D}-\mathbf{M}(\theta)\big)^{\!\top}\mathbf{C}^{-1}
\big(\mathbf{D}-\mathbf{M}(\theta)\big),
\label{eq:lnL_CMB}
\end{equation}
with covariance $\mathbf{C}$ provided by the likelihood package.
Model selection uses $\Delta\chi^2=-2\Delta\ln\mathcal{L}$ against the $\Lambda$CDM baseline at matched
nuisance treatment.


\paragraph{Implementation notes (CLASS/CAMB + Cobaya/MontePython).}
A minimal reproducible pipeline is: (i) implement the discrete-$k$ sums and any $\phig$-indexed feature parameters in a closed-geometry Boltzmann solver
(CLASS/CAMB)~\cite{Blas2011CLASS,Lewis2000CAMB}; (ii) couple to an MCMC / nested-sampling front-end
(Cobaya or MontePython)~\cite{TorradoLewis2021Cobaya,Audren2013MontePython} using the official Planck likelihoods;
(iii) publish the full chain set and posterior plots with a standard analyzer (e.g.\ GetDist)~\cite{Lewis2019GetDist}.

\paragraph{Stochastic gravitational-wave background.}
Given a prediction $\Omega_{\rm GW}(f;\theta)$ (Sec.~\ref{sec:sgwb}), compare to LVK~\cite{LVK2021ISGWB} and PTA constraints via
their published likelihoods. A simple summary statistic is
\begin{equation}
\chi^2_{\rm GW}(\theta)=\sum_i \frac{\big(\Omega_{\rm GW}(f_i;\theta)-\Omega^{\rm obs}_i\big)^2}{\sigma_i^2},
\label{eq:chi2_GW}
\end{equation}
for an appropriate set of frequency bins $(f_i,\Omega^{\rm obs}_i,\sigma_i)$, while the definitive test
uses the collaboration likelihoods.

\paragraph{Non-negotiable deliverables.}
A complete submission must therefore include: (i) a reproducible code path generating $C_\ell$ and
$\Omega_{\rm GW}(f)$ from the manuscript equations, (ii) a parameter table with priors, and
(iii) a $\Delta\chi^2$ (or Bayes factor) report against standard baselines.


\section{Stochastic gravitational-wave background}
\label{sec:sgwb}
\subsection{Discrete tensor spectrum and frequencies}
Tensor harmonics on $\Sthree$ form a discrete set:
\begin{equation}
\Delta_{\Sthree}T^{(n)}_{ij}=-\frac{\lambda^{(T)}_n}{R_K^2}\,T^{(n)}_{ij},\qquad n\ge 3,
\label{eq:S3_tensor_eigs}
\end{equation}
with $k^{(T)}_n\sim \sqrt{\lambda^{(T)}_n}/R_K$ and $f_n=\frac{c}{2\pi}\frac{k^{(T)}_n}{a_0}$.

\subsection{\texorpdfstring{Definition and qualitative signature of $\Omega_{\rm GW}(f)$}{Definition and qualitative signature of OmegaGW(f)}}
\begin{equation}
\Omega_{\rm GW}(f)\equiv \frac{1}{\rho_c}\frac{\dd\rho_{\rm GW}}{\dd\ln f},
\qquad \rho_c=\frac{3H_0^2c^2}{8\pi G}.
\label{eq:OmegaGW_def}
\end{equation}
Discreteness yields (schematically)
\begin{equation}
\Omega_{\rm GW}(f)\sim \sum_{n\ge 3}\mathcal P_T(k^{(T)}_n)\,\mathcal T_T^2(k^{(T)}_n)\,\delta(f-f_n),
\label{eq:OmegaGW_discrete}
\end{equation}
implying an IR cutoff and possible step/comb features set by $1/(a_0R_K)$.

Equation~\eqref{eq:OmegaGW_discrete} is schematic: in any realistic analysis the \(\delta\)-comb is broadened
by finite observation time, redshift smearing, and the bandwidth of estimators used in PTA/LVK searches.
Moreover, for curvature radii compatible with current near-flat constraints, the mode spacing in the PTA/LVK bands is
effectively continuous: one has $\Delta f/f\sim 1/n$ with $n\sim 2\pi f\,a_0R_K/c\gg 1$.
Current PTA detections are well-fit by a smooth power-law spectrum consistent with astrophysical SMBHB
populations~\cite{NANOGrav15yr2023,IPTADR2_2022}; therefore topology-induced discreteness must either lie in a
marginally allowed regime (very large \(R_K\)) or appear only as weak modulations on top of a smooth background.


\subsection{Minimal SGWB verification protocol}
\begin{enumerate}
\item Fix $R_K$ and specify $\mathcal P_T(k)$.
\item Compute $(k^{(T)}_n,f_n)$.
\item Build broadened $\Omega_{\rm GW}(f)$ prediction.
\item Compare to PTA and ground-based constraints~\cite{NANOGrav15yr2023,IPTADR2_2022,LVK2021ISGWB}; forecast detectability of topology-induced features.
\end{enumerate}

\section{Conclusion}
Diffeomorphism invariance yields Noether-II identities rather than global Noether-I charges. Associated covariant
charges are boundary charges and vanish on compact boundaryless $\Sig$ (notably $\Sthree$). Embedding-level
modifications are variationally well-defined either as smooth $\nu(t)$ dependence or as discrete floor $N(t)$ with
explicit junction consistency. Compact topology implies discrete mode spectra, giving falsifiable signatures in low-$\ell$
CMB and infrared SGWB, testable through likelihood analyses.

Present data already impose strong constraints on both spatial curvature and compact topology. In particular,
Planck cosmological parameter analyses tightly constrain curvature around \(\Omega_k\approx 0\)~\cite{Planck2018Cosmo},
while dedicated ``matched circles'' and related searches constrain the size of many compact fundamental domains to be
comparable to or larger than the surface of last scattering~\cite{Planck2015Topology,Saha2024LensSpaces,FutureTopologySearches}.
Accordingly, the discrete-mode signatures emphasized here are expected to appear only in the marginally allowed tail of
parameter space (very large \(R_K\) and/or nearly-horizon fundamental domains), with upcoming CMB and PTA datasets providing
decisive tests.

\paragraph{Invariant--constraint map.}
To make the falsifiable content explicit, we summarize the chain \emph{invariant \(\Rightarrow\) observable \(\Rightarrow\) present constraint \(\Rightarrow\) likelihood module}:
\begin{enumerate}
\item \textbf{Diffeomorphism invariance \(\Rightarrow\) Noether-II identities and boundary charges.}
In generally covariant theories the on-shell diffeomorphism charge is a boundary integral~\cite{Noether1918,IyerWald1994}.
Observable: vanishing global diffeomorphism charges on compact boundaryless slices (\(\partial\Sigma=\varnothing\)).
Constraint/module: the analysis is implemented at the level of covariant identities and requires no asymptotic boundary data.

\item \textbf{Spherical space form sector \(\widetilde\Sigma\simeq S^3\) \(\Rightarrow\) discrete Laplacian spectra.}
For \(\Sigma\simeq S^3/\Gamma\) the scalar/tensor mode spectra are discrete~\cite{WolfConstantCurvature,Scott3Manifolds}.
Observable: low-\(\ell\) discreteness/cutoff in \(C_\ell\).
Constraint/module: tested against Planck likelihood and curvature constraints~\cite{Planck2018Lik,Planck2018Cosmo} and topology searches~\cite{Planck2015Topology,Saha2024LensSpaces,FutureTopologySearches}.

\item \textbf{Invariant refinement bookkeeping \(\Rightarrow\) unique scale ladder \(\Rightarrow\) \(\varphi\).}
Requiring unambiguous composition of refinement steps fixes the Zeckendorf--Lekkerkerker encoding and hence the Fibonacci ladder, with \(\varphi\) as the unique positive fixed point of the induced ratio map~\cite{Lekkerkerker1951,Keller1972}.
Observable: step/log-periodic features tied to \(N(t)\) transitions.
Constraint/module: junction conditions enforce absence of spurious distributional sources (Sec.~IV).

\item \textbf{Compactness \(+\) tensor discreteness \(\Rightarrow\) IR structure in \(\Omega_{\rm GW}(f)\).}
Observable: broadened IR features (rather than idealized delta-combs) in the PTA/LVK bands.
Constraint/module: compared to PTA and ground-based bounds~\cite{NANOGrav15yr2023,IPTADR2_2022,LVK2021ISGWB}.
\end{enumerate}

\section{Verification checklist}
\begin{enumerate}
\item Implement closed-geometry mode functions; compute $C_\ell$ from \eqref{eq:Cl_discrete_sum}.
\item Fit Planck low-$\ell$ likelihood; constrain $R_K$ and topology parameters.
\item Implement tensor discreteness; compare $\Omega_{\rm GW}(f)$ to PTA and ground-based bounds.
\item Report likelihood metrics ($\chi^2$, Bayes factors), not qualitative claims.
\end{enumerate}



\appendix

\section{\texorpdfstring{From invariant assumptions to $S^{3}$: hard vs.\ soft inevitability}{From invariant assumptions to S3: hard vs soft inevitability}}
\label{app:hard_soft_s3}

\subsection{Invariant, sector, and test}
\label{app:invariant_sector_test}
In this paper, ``invariant'' means coordinate-free and gauge-independent. Diffeomorphism invariance in GR
yields Noether-II identities that constrain the local form of the field equations and the boundary-term
structure, but it does not uniquely fix global topology. A global assumption such as compactness and
boundarylessness ($\partial\Sig=\varnothing$) specifies a topological sector of configuration space, while
additional topological/geometric conditions can make the spherical sector mathematically unavoidable.

Hard inevitability corresponds to the topological invariant $\pi_1(\Sig)=0$, forcing $\Sig\cong \Sthree$
(Theorem~\ref{thm:hard_s3_main}). Soft inevitability corresponds to the geometric invariant $\mathrm{Ric}>0$,
forcing spherical space forms $\Sig\cong \Sthree/\Gamma$ with $\widetilde{\Sig}\cong \Sthree$
(Corollary~\ref{cor:soft_s3_main}). In practice, present data typically constrain the size and group label
$\Gamma$ rather than proving trivial $\Gamma$; hence we implement explicit likelihood comparison within the
spherical sector.

\subsection{Likelihood selection within the spherical sector}
\label{app:likelihood_spherical_sector}
Within $\Sig\cong \Sthree/\Gamma$, topology enters observables through the Laplacian eigenmodes and the induced
correlation structure of perturbations (and their $\phig$-refined discretization). This enables model
comparison via Bayesian evidence and Bayes factors:
\begin{equation}
\mathcal{Z}_{\Gamma}=\int d\theta\, \mathcal{L}(D\mid \theta,\Gamma)\,\pi(\theta\mid \Gamma),
\qquad
B_{\Gamma_1\Gamma_2}=\frac{\mathcal{Z}_{\Gamma_1}}{\mathcal{Z}_{\Gamma_2}}.
\end{equation}
Here $\Gamma$ labels the (finite) spherical space-form group and $\theta$ collects cosmological and nuisance
parameters. This is the operational implementation of the invariant $\to$ constraint $\to$ observable $\to$
likelihood map used throughout the manuscript.


\section{Invariant sieve to \texorpdfstring{$S^{3}$}{S3}: properties, theorems, and experimental handles}
\label{app:invariant_sieve}
The goal of this appendix is to make explicit the ``invariant sieve'' logic used in the main text.
We start from a broad candidate class of spatial topologies $\Sig$ (closed, smooth 3-manifolds), and we
collect coordinate-free and gauge-independent inputs (``invariants'') that translate into topological or
geometric constraints. The admissible class is then the intersection
\begin{equation}
\mathfrak{T}_{\rm adm} \;=\; \bigcap_{i}\,\mathcal{C}_{i},
\qquad
\mathcal{C}_{i}\subset \mathfrak{T}_{0}\;\;\text{(candidate class)},
\end{equation}
where each $\mathcal{C}_{i}$ is induced by an invariant statement and/or a data-driven constraint.
The sieve is strongest when (i) the mathematical implication ``invariant $\Rightarrow$ constraint'' is
theorem-level, and (ii) the associated observable is actually testable by a concrete likelihood analysis.

\vspace{0.5em}
\noindent\textbf{What can and cannot be ``uniquely forced'' by experiment.}
Current cosmological topology searches primarily provide bounds on the size of the fundamental domain and
exclude specific nontrivial identifications up to some scale, rather than proving uniqueness of
$\pi_{1}(\Sig)$ or of the group label $\Gamma$ \cite{CornishSpergelStarkman1998,Planck2015Topology}.
Accordingly, the sieve has two layers:
(a) a theorem-driven layer (hard/soft inevitability) that can uniquely single out $\Sthree$ under explicit
topological/geometric invariants, and
(b) a likelihood layer that constrains (or rules out) candidate quotients $\Sthree/\Gamma$.

\begin{table*}[t]
\centering
\renewcommand{\arraystretch}{1.15}
\begin{tabular}{p{0.26\textwidth} p{0.38\textwidth} p{0.32\textwidth}}
\hline\hline
\textbf{Property of $S^{3}$ (invariant)} & \textbf{Theorem-level implication} & \textbf{Experimental / observational handle} \\
\hline
Closed and boundaryless ($\partial\Sig=\varnothing$) & In GR, diffeo-charges are boundary charges; for
$\partial\Sig=\varnothing$ the corresponding global boundary integrals vanish (Sec.~\ref{sec:explicit-generalized-noether}).
This fixes the correct meaning of ``conservation'' in the compact sector. & Not directly measurable as a
binary fact; motivates focusing on topology-dependent correlation structure rather than global energy.
\\
\hline
Simply connected ($\pi_{1}(\Sthree)=0$) & \emph{Hard inevitability:} for a closed 3-manifold, $\pi_{1}(\Sig)=0$
implies $\Sig\cong \Sthree$ (Poincar\'e/Perelman; used as Theorem~\ref{thm:hard_s3_main}). & CMB topology
searches (matched circles and related correlation tests) constrain nontrivial identifications and hence
possible nontrivial $\pi_{1}(\Sig)$ at scales comparable to the LSS radius \cite{CornishSpergelStarkman1998,Planck2015Topology}.
\\
\hline
Spherical space-form sector (constant $K>0$) & If $\Sig$ admits a metric of constant positive sectional
curvature then $\Sig\cong \Sthree/\Gamma$ with $\widetilde\Sig\cong \Sthree$ (spherical space-form
classification; used in Corollary~\ref{cor:soft_s3_main} as sector target). & CMB topology constraints
restrict specific $\Gamma$ and the fundamental-domain size; likelihood analysis compares $\Gamma$-labels via
mode spectra and correlation structure (Appendix~\ref{app:likelihood_spherical_sector}).
\\
\hline
Positive Ricci curvature route ($\mathrm{Ric}>0$) & \emph{Soft inevitability:} $\mathrm{Ric}>0$ on a closed
3-manifold forces a spherical space form $\Sthree/\Gamma$ (Appendix~\ref{app:hard_soft_s3}). & Not directly
measured as a geometric invariant of spatial slices; implemented as a physically motivated sector prior.
\\
\hline
Discrete Laplacian spectrum on $\Sthree$ & Compactness implies a discrete spectrum; on $\Sthree$ the
eigenmodes are hyperspherical harmonics, enabling replacement of $\int dk$ by a mode sum (Sec.~\ref{sec:topology_status}, Eq.~\eqref{eq:Cl_discrete_sum}). & Low-$\ell$ CMB power and correlation anomalies are sensitive to IR mode content;
topology searches and Planck likelihood constrain topological parameters and curvature scale \cite{Planck2015Topology}.
\\
\hline
Spin structure / global fermions & Any oriented 3-manifold is spin (vanishing $w_2$); $\Sthree$ is spin and
admits globally defined spinors, consistent with including fermionic quantum fields. & Laboratory quantum
experiments show that topological information can be encoded in gauge-invariant phases and discrete
indices (e.g. Aharonov--Bohm and quantum Hall invariants)\cite{AharonovBohm1959,Klitzing1980,TKNN1982}.
This supports ``topology can be observable'' but does not by itself fix cosmic $\Sig$.
\\
\hline\hline
\end{tabular}
\caption{\label{tab:invariant_sieve}
Invariant sieve: representative topological/geometric properties of $S^{3}$, theorems that map them to
restrictions on admissible $\Sig$, and the corresponding experimental/observational handles.
Theorem-level inputs (hard/soft) can single out the spherical sector, while data-driven searches constrain
the quotient label $\Gamma$ and the fundamental-domain size within that sector.
}
\end{table*}

\noindent\textbf{Sieve conclusion.}
The combined effect of (i) Noether-consistency in the compact sector (boundary-charge structure),
(ii) theorem-level hard/soft routes to the spherical sector, and
(iii) explicit likelihood comparison over $\Gamma$ via topology-dependent mode spectra,
is to reduce the admissible topology class to the spherical sector $\Sthree/\Gamma$ and to treat $\Sthree$
as the minimal representative. A mathematically unique identification $\Sig\cong \Sthree$ is obtained once
the hard invariant $\pi_{1}(\Sig)=0$ is established or adopted as a required constraint; observationally,
this corresponds to progressively excluding nontrivial identifications through topology searches and
likelihood-based model comparison.

% --- Flush any pending floats (especially \texttt{table*}) before starting the bibliography.
% RevTeX can otherwise place deferred wide floats inside the reference section.
\clearpage


\begin{thebibliography}{99}

\bibitem{Noether1918}
E.~Noether,
``Invariante Variationsprobleme,''
Nachr.\ Ges.\ Wiss.\ G{\"o}ttingen, Math.-Phys.\ Kl.\ \textbf{1918}, 235--257 (1918).

\bibitem{BradingBrown2003}
K.~Brading and H.~R.~Brown,
``Noether's Theorems and Gauge Symmetries,''
Stud.\ Hist.\ Philos.\ Mod.\ Phys.\ \textbf{34}, 509--538 (2003).

\bibitem{LeeWald1990}
J.~Lee and R.~M.~Wald,
``Local symmetries and constraints,''
J.\ Math.\ Phys.\ \textbf{31}, 725--743 (1990).

\bibitem{WaldBook}
R.~M.~Wald,
\emph{General Relativity}
(University of Chicago Press, Chicago, 1984).

\bibitem{IyerWald1994}
V.~Iyer and R.~M.~Wald,
``Some properties of Noether charge and a proposal for dynamical black hole entropy,''
Phys.\ Rev.\ D \textbf{50}, 846 (1994),
arXiv:gr-qc/9403028.

\bibitem{PreeceBatenin2025QCNoether}
A.~Preece and B.~Batenin,
``A Generalized Noether's Theorem in Quantum Cosmology,''
preprint (2025).

\bibitem{PreeceBatenin2025MFNoether}
A.~Preece and B.~Batenin,
``A Generalized Noether's Theorem for Multifractal Cosmology,''
preprint (2025).

\bibitem{CornishSpergelStarkman1998}
N.~J.~Cornish, D.~N.~Spergel, and G.~D.~Starkman,
``Circles in the Sky: Finding Topology with the Microwave Background Radiation,''
Phys.\ Rev.\ D \textbf{57}, 5982 (1998),
arXiv:astro-ph/9801212.

\bibitem{Planck2015Topology}
Planck Collaboration,
P.~A.~R.~Ade \textit{et al.},
``Planck 2015 results. XVIII. Background geometry and topology of the Universe,''
Astron.\ Astrophys.\ \textbf{594}, A18 (2016),
arXiv:1502.01593.

\bibitem{Saha2024LensSpaces}
COMPACT Collaboration,
S.~Saha \textit{et al.},
``Cosmic topology. Part Ic. Limits on lens spaces from circle searches,''
preprint (2024),
arXiv:2409.02226.

\bibitem{FutureTopologySearches}
COMPACT Collaboration,
Y.~Akrami \textit{et al.},
``Promise of Future Searches for Cosmic Topology,''
Phys.\ Rev.\ Lett.\ \textbf{132}, 171501 (2024),
arXiv:2210.11426.

\bibitem{WolfConstantCurvature}
J.~A.~Wolf,
\emph{Spaces of Constant Curvature},
6th ed. (AMS Chelsea Publishing, Providence, RI, 2011).

\bibitem{Scott3Manifolds}
P.~Scott,
``The geometries of 3-manifolds,''
Bull.\ Lond.\ Math.\ Soc.\ \textbf{15}, 401--487 (1983).


\bibitem{Hamilton1982PositiveRicci}
R.~S.~Hamilton,
``Three-manifolds with positive Ricci curvature,''
J.\ Diff.\ Geom.\ \textbf{17}, 255--306 (1982).

\bibitem{Perelman2002Entropy}
G.~Perelman,
``The entropy formula for the Ricci flow and its geometric applications,''
arXiv:math/0211159.

\bibitem{Perelman2003Surgery}
G.~Perelman,
``Ricci flow with surgery on three-manifolds,''
arXiv:math/0303109.

\bibitem{Perelman2003Extinction}
G.~Perelman,
``Finite extinction time for the solutions to the Ricci flow on certain three-manifolds,''
arXiv:math/0307245.

\bibitem{MorganTian2007RicciFlowPoincare}
J.~Morgan and G.~Tian,
\emph{Ricci Flow and the Poincar\'e Conjecture}
(Clay Mathematics Monographs, Vol.~3, American Mathematical Society, Providence, RI, 2007).

\bibitem{KleinerLott2008NotesPerelman}
B.~Kleiner and J.~Lott,
``Notes on Perelman{'}s papers,''
Geom.\ Topol.\ \textbf{12}, 2587--2858 (2008),
arXiv:math/0605667.

\bibitem{Planck2018Lik}
Planck Collaboration,
N.~Aghanim \textit{et al.},
``Planck 2018 results. V. CMB power spectra and likelihoods,''
Astron.\ Astrophys.\ \textbf{641}, A5 (2020),
arXiv:1907.12875.

\bibitem{Planck2018Cosmo}
Planck Collaboration,
N.~Aghanim \textit{et al.},
``Planck 2018 results. VI. Cosmological parameters,''
Astron.\ Astrophys.\ \textbf{641}, A6 (2020),
arXiv:1807.06209.

\bibitem{Blas2011CLASS}
D.~Blas, J.~Lesgourgues, and T.~Tram,
``The Cosmic Linear Anisotropy Solving System (CLASS). Part II: Approximation schemes,''
JCAP \textbf{07}, 034 (2011),
arXiv:1104.2932.

\bibitem{Lewis2000CAMB}
A.~Lewis, A.~Challinor, and A.~Lasenby,
``Efficient computation of CMB anisotropies in closed FRW models,''
Astrophys.\ J.\ \textbf{538}, 473 (2000),
arXiv:astro-ph/9911177.

\bibitem{TorradoLewis2021Cobaya}
J.~Torrado and A.~Lewis,
``Cobaya: Code for Bayesian analysis of hierarchical physical models,''
JCAP \textbf{05}, 057 (2021),
arXiv:2005.05290.

\bibitem{Audren2013MontePython}
B.~Audren, J.~Lesgourgues, K.~Benabed, and S.~Prunet,
``Conservative constraints on early cosmology with Monte Python,''
JCAP \textbf{02}, 001 (2013),
arXiv:1210.7183.

\bibitem{Lewis2019GetDist}
A.~Lewis,
``GetDist: a Python package for analysing Monte Carlo samples,''
arXiv:1910.13970 [astro-ph.IM] (2019).

\bibitem{LVK2021ISGWB}
R.~Abbott \textit{et al.} (LIGO Scientific Collaboration and Virgo Collaboration),
``Upper Limits on the Isotropic Gravitational-Wave Background from Advanced LIGO and Advanced Virgo's Third Observing Run,''
Phys.\ Rev.\ D \textbf{104}, 022004 (2021),
arXiv:2101.12130.

\bibitem{NANOGrav15yr2023}
G.~Agazie \textit{et al.} (NANOGrav Collaboration),
``The NANOGrav 15-year Data Set: Evidence for a Gravitational-Wave Background,''
Astrophys.\ J.\ Lett.\ \textbf{951}, L18 (2023),
arXiv:2306.16213.

\bibitem{IPTADR2_2022}
J.~Antoniadis \textit{et al.} (International Pulsar Timing Array Collaboration),
``The International Pulsar Timing Array second data release: Search for an isotropic gravitational wave background,''
Mon.\ Not.\ R.\ Astron.\ Soc.\ \textbf{510}, 4873 (2022),
arXiv:2201.03980.

\bibitem{AharonovBohm1959}
Y.~Aharonov and D.~Bohm,
``Significance of electromagnetic potentials in the quantum theory,''
Phys.\ Rev.\ \textbf{115}, 485 (1959).

\bibitem{Klitzing1980}
K.~v.~Klitzing, G.~Dorda, and M.~Pepper,
``New method for high-accuracy determination of the fine-structure constant based on quantized Hall resistance,''
Phys.\ Rev.\ Lett.\ \textbf{45}, 494 (1980).

\bibitem{TKNN1982}
D.~J.~Thouless, M.~Kohmoto, M.~P.~Nightingale, and M.~den~Nijs,
``Quantized Hall Conductance in a Two-Dimensional Periodic Potential,''
Phys.\ Rev.\ Lett.\ \textbf{49}, 405 (1982).


\bibitem{Lekkerkerker1951}
C.~G.~Lekkerkerker,
``Voorstelling van natuurlijke getallen door een som van getallen van Fibonacci,''
Simon Stevin \textbf{29} (1951--1952), 190--195.

\bibitem{Keller1972}
T.~J.~Keller,
``Generalizations of Zeckendorf's Theorem,''
Fibonacci Quarterly \textbf{10} (1972), no.~1, 95--102.

\bibitem{Chang2020WeakConverse}
S.~Chang,
``The weak converse of Zeckendorf's Theorem,''
arXiv:2007.12169 [math.NT] (2020).

\end{thebibliography}


\end{document}
